\section*{Chapter \ref{ch:qm:floating-point-and-numerical-linear-algebra} --- Floating Point and Numerical Linear Algebra}

\begin{solution}{exo:qm:floating-point-and-numerical-linear-algebra:1}
Unit roundoff $u = 2^{-53} \approx 1.11 \times 10^{-16}$, \omterm{def:qm:floating-point-and-numerical-linear-algebra:fp}{machine epsilon} $2^{-52} \approx 2.22 \times 10^{-16}$; $53\log_{10}2 \approx 15.95$, so fifteen to sixteen significant decimal digits.
\end{solution}

\begin{solution}{exo:qm:floating-point-and-numerical-linear-algebra:2}
Half a hundredth per recalculation: $0.005 \times 10\,000 \times 250 = 12\,500$ points a year, downward, before the compounding with the index's own moves.
\end{solution}

\begin{solution}{exo:qm:floating-point-and-numerical-linear-algebra:3}
$10^8 + 1$ needs 27 bits and is representable, so both operations are exact. Above $2^{53} \approx 9.0 \times 10^{15}$ the spacing of doubles is 2, so $10^{16} + 1$ is a tie, rounded to the
neighbour with an even significand, which is $10^{16}$: the result is 0.
\end{solution}

\begin{solution}{exo:qm:floating-point-and-numerical-linear-algebra:4}
Each of $\sum x_i^2$ and $(\sum x_i)^2/n$ is about $nc^2$ and carries a rounding error of at least $u\,nc^2$ (more with recursive summation); their difference, divided by $n - 1$, has an absolute
error of order $uc^2$. With $c = 10^4$: $1.1 \times 10^{-16} \times 10^8 \approx 10^{-8}$, or about $10^{-7}$ relative to a variance of 0.08; the summation error pushes it up, and the tutorial
measures $8.0 \times 10^{-7}$.
\end{solution}

\begin{solution}{exo:qm:floating-point-and-numerical-linear-algebra:5}
The normal equations work with $\kappa^2 = 10^{10}$ and can lose about ten of sixteen digits; QR, with $\kappa = 10^5$, about five.
\end{solution}

\begin{solution}{exo:qm:floating-point-and-numerical-linear-algebra:6}
All the $x_i - m$ are at most 0, so each exponential is at most 1 and the largest is exactly 1: the sum lies in $[1, n]$ and its logarithm cannot overflow or be $-\infty$. $\log(e^{1000} + e^{1000}) =
1000 + \ln2 = 1000.693$, where the direct formula overflows to infinity.
\end{solution}

\begin{solution}{exo:qm:floating-point-and-numerical-linear-algebra:7}
\texttt{summation\_errors()}: at $n = 10^6$ the errors against the exact rational sum (47\,878\,411.409) are $2.5 \times 10^{-6}$ (naive), $2.8 \times 10^{-8}$ (pairwise) and $2.2 \times 10^{-9}$
(Neumaier), the last below one \omterm{def:qm:floating-point-and-numerical-linear-algebra:fp}{unit in the last place} ($7.5 \times 10^{-9}$).
\end{solution}

\begin{solution}{exo:qm:floating-point-and-numerical-linear-algebra:8}
A difference in the tenth digit is what different operation orders, \omterm{def:qm:floating-point-and-numerical-linear-algebra:fma}{fused multiply-add} contraction or different \texttt{exp} and \texttt{log} implementations produce on a moderately conditioned
computation; it is not evidence of a bug. Evidence would be a difference larger than the \omterm{def:qm:floating-point-and-numerical-linear-algebra:cond}{condition number} times the unit roundoff, or a difference on reference inputs where both engines are meant
to perform identical operations. Either fix the order and flags and assert bits, or compare with a tolerance justified by the conditioning.
\end{solution}

\begin{solution}{pb:qm:floating-point-and-numerical-linear-algebra:1}
\textbf{1.} If the digits beyond the third decimal are uniformly distributed, the discarded part is uniform on $[0, 0.001)$: 0.0005 on average, always a loss.
\textbf{2.} $0.0005 \times 2\,400 \times 480 = 576$ points, against $1098.892 - 524.811 = 574.081$.
\textbf{3.} $0.0005 \times 3\,000 \times 480 = 720$ points.
\textbf{4.} Rounding errors are symmetric around zero: no drift, only noise of standard deviation about $0.001\sqrt{n/12} \approx 0.3$; the simulated rounded index ends at 1098.57 against 1098.892.
\textbf{5.} At 422.23, a loss of 676.7. A point lost at step $k$ is carried forward by the index's growth from $k$ to the end, so the loss is $0.0005\sum_kE_n/E_k$; the simulated exact index
spends most of its time below its final value (between 803 and 1183), and the average ratio is 1.175, so $576 \times 1.175 = 676.7$.
\textbf{6.} $2.5 \times 10^{-6}$, $2.8 \times 10^{-8}$ and $2.2 \times 10^{-9}$.
\textbf{7.} Relative errors $8.0 \times 10^{-7}$, 1.3\% and 1060 (86.84 for a true 0.0819).
\textbf{8.} They subtract the mean before squaring, so they square small deviations; their errors are of order $u$ times the spread, not $u$ times the offset squared (about $10^{-12}$ at most for two
passes, $1.6 \times 10^{-8}$ for Welford at an offset of $10^8$).
\textbf{9.} Relative errors $5.3 \times 10^{-3}$ and $1.2 \times 10^{-10}$: the normal equations keep about two digits, QR about ten (seven lost, as $\kappa = 10^7$ predicts).
\textbf{10.} At the sixth pivot (index 5), which is $-0.14$: the pairwise-estimated matrix is not positive semidefinite (smallest eigenvalue $-1.42$), so by
\cref{prop:qm:floating-point-and-numerical-linear-algebra:chol} the factorisation must fail.
\textbf{11.} The bits of eight results (the naive, pairwise and Neumaier sums, Welford's mean and variance, a log-sum-exp, the first and last components of a tridiagonal solve), asserted as the same
hexadecimal patterns in the Python, C++20 and Rust tests.
\textbf{12.} Two roundings give exactly 0; one rounding gives $2^{-54} \approx 5.6 \times 10^{-17}$, the representation error of $0.1$ times 10.
\textbf{13.} \texttt{-ffp-contract=fast}, GCC's default outside a standards-compliant C mode, once the target (\texttt{-march}) has a \omterm{def:qm:floating-point-and-numerical-linear-algebra:fma}{fused multiply-add}; also \texttt{-ffast-math}.
\textbf{14.} 1\,909 without and 102 with the \omterm{def:qm:floating-point-and-numerical-linear-algebra:cg}{preconditioner}.
\textbf{15.} Rounding destroys the conjugacy of the directions, so the finite-termination argument fails; what survives is the $\sqrt\kappa$ rate ($\sqrt{1.1 \times 10^5} \approx 330$), whose bound
predicts about 3\,100 iterations for a factor $10^8$; the observed 1\,909 is within it.
\textbf{16.} Recompute the index from scratch, from constituent prices and the divisor, once a day and compare with the running value: a gap growing by about a point a day is visible in a week.
\textbf{17.} Input hashes, code and library versions, compiler and flags, thread count and reduction order, seeds, and the bit patterns of key intermediate aggregates.
\textbf{18.} When numbers are reconciled across systems or over time (limits, regulatory reports, regression tests of a backtest), where an unexplained difference costs investigation time; the cost is
fixed-order reductions and no fast-math.
\textbf{19.} Named result: \emph{the index that lost half its value}: truncation removes 0.0005 per recalculation on average, $0.0005 \times 2\,400 \times 480 = 576$ points, against a documented
shortfall of 574.081 (524.811 against 1098.892).
\textbf{20.} It rounds it at every operation, with relative error at most $2^{-53}$, and what the algorithm and the conditioning do with those roundings decides the answer.
\end{solution}

\begin{solution}{iq:qm:floating-point-and-numerical-linear-algebra:1}
Neither 0.1 nor 0.2 nor 0.3 is representable in binary; the rounded sum of the first two is one \omterm{def:qm:floating-point-and-numerical-linear-algebra:fp}{unit in the last place} above the rounded 0.3. Compare with a tolerance mixing relative and absolute
parts, $|a - b| \le \max(r\max(|a|, |b|), t)$, chosen from the computation's conditioning; or work in integers (ticks, cents) where exactness is required.
\iqlookfor{Binary representation, and a relative-plus-absolute tolerance or integer units.}
\end{solution}

\begin{solution}{iq:qm:floating-point-and-numerical-linear-algebra:2}
Different operation order (vectorised or parallel reductions, thread counts), \omterm{def:qm:floating-point-and-numerical-linear-algebra:fma}{fused multiply-add} contraction, fast-math flags, different transcendental implementations (\texttt{exp}, \texttt{log},
\texttt{pow} are not correctly rounded), extended precision in intermediate registers, different parsing of decimal inputs, and different library algorithms (a BLAS versus a hand loop).
\iqlookfor{Non-associativity and at least three concrete sources.}
\end{solution}

\begin{solution}{iq:qm:floating-point-and-numerical-linear-algebra:3}
Welford's update: $\delta = x - m$, $m \leftarrow m + \delta/n$, $M_2 \leftarrow M_2 + \delta(x - m)$. The sum of squares minus the squared sum cancels catastrophically when the mean is large
relative to the spread, as prices are; at prices near $10^8$ it gives nonsense.
\iqlookfor{Welford's recurrence and the cancellation argument.}
\end{solution}

\begin{solution}{iq:qm:floating-point-and-numerical-linear-algebra:4}
The normal equations form $X^\top X$, whose \omterm{def:qm:floating-point-and-numerical-linear-algebra:cond}{condition number} is $\kappa(X)^2$, and lose twice as many digits; QR works on $X$ directly and is backward stable. At $\kappa = 10^7$ the chapter's
coefficients are wrong in the third digit by the normal equations and right to ten by QR.
\iqlookfor{Squared \omterm{def:qm:floating-point-and-numerical-linear-algebra:cond}{condition number}.}
\end{solution}

\begin{solution}{iq:qm:floating-point-and-numerical-linear-algebra:5}
First ask why: a matrix estimated pairwise or with missing data can be indefinite, and a sample covariance with fewer observations than assets is singular. Look at the failing index and the
eigenvalues. Then repair with the \omterm{def:qm:convex-optimisation:ncm}{nearest correlation matrix}, shrinkage, or a \omterm{def:qm:covariance-estimation-and-random-matrices:factor}{factor model}; a small jitter is acceptable only for a matrix that is positive semidefinite up to rounding. Never
jitter silently.
\iqlookfor{Diagnosis before repair, and a principled repair.}
\end{solution}

\begin{solution}{iq:qm:floating-point-and-numerical-linear-algebra:6}
Fix the decomposition of the data into blocks independently of the thread count, sum each block in a fixed order (with compensation for accuracy), and combine the partial sums in a fixed
tree; threads then only schedule blocks. Alternatives: exact accumulation (a long accumulator, or integer fixed point), or a pre-rounding scheme that makes addition exact.
\iqlookfor{Order independent of the thread count.}
\end{solution}
